% ======================================================================
\section{Results}\label{sec:results}

All numbers in this section are raw counts from the scorer (\fp{benchmark/gold/score_per_sound.py}) or values derived
from those counts by Eq.~\ref{eq:cost}--\ref{eq:costw}; derived values are marked. \Dp{} counts are reproduced
exactly by the Decision Inspector export (\fp{docs/inspector2/data_parity.json}, \texttt{"exact": true}).

\subsection{Headline}\label{sec:headline}
\begin{table}[h]
\centering\small
\caption{Main result on the merged sets. Recall = hits / needed; precision = hits / (hits + wrong) (derived). ``Show
nothing'' is the cost of a system that never shows a picture, $4\cdot\text{needed}/\text{clips}$ (derived).}
\label{tab:headline}
\resizebox{\linewidth}{!}{%
\begin{tabular}{lrrrrrrrr}
\toprule
 & \multicolumn{4}{c}{merged DEV (71 clips, 58 needed)} & \multicolumn{4}{c}{merged TEST (88 clips, 65 needed)} \\
\cmidrule(lr){2-5}\cmidrule(lr){6-9}
system & hits & wrong (v/c/p) & cost & R / P & hits & wrong (v/c/p) & cost & R / P \\
\midrule
show nothing & 0 & 0 & 3.268 & -- & 0 & 0 & 2.955 & -- \\
old pipeline (B0r) & 18 & 51 (5/37/9) & 3.690 & .31 / .26 & 21 & 40 (2/31/7) & 2.909 & .32 / .34 \\
B (SHIP8+MD3) & 29 & 18 (6/9/3) & 2.141 & .50 / .62 & 22 & 26 (4/17/5) & 2.545 & .34 / .46 \\
D (text readout) & 29 & 14 (6/6/2) & 2.028 & .50 / .67 & 24 & 23 (4/14/5) & 2.386 & .37 / .51 \\
\textbf{\Dp{} (frozen)} & \textbf{29} & \textbf{15 (6/7/2)} & \textbf{2.056} & \textbf{.50 / .66} & \textbf{24} & \textbf{24 (4/15/5)} & \textbf{2.409} & \textbf{.37 / .50} \\
\bottomrule
\end{tabular}}

\smallskip
\footnotesize v / c / p = wrong pictures whose source is on screen / that belong to another sound / with nothing
there. Sources: \fp{docs/review/detector_final_2026-10-02.md}; \fp{docs/review/improvements_log_2026-09-30.md} (B0r);
\fp{docs/prereg_round13_detector_push.md} lines 3885--3888, 4219--4243, 4536--4544, 4631--4633.
\end{table}

\paragraph{Reading the table.}
\begin{itemize}
  \item On the development set \Dp{} finds half of the needed sounds and two out of three of its pictures are right.
        On the test set it finds 37\,\% and half of its pictures are right.
  \item The old pipeline was \emph{worse than showing nothing} on the development set (3.690 vs 3.268; still 3.521
        with a 3-s repeat merge, Section~\ref{sec:mergegap}) and only
        slightly better on the test set (2.909 vs 2.955). \Dp{} is clearly cheaper than showing nothing on both
        (2.056 vs 3.268; 2.409 vs 2.955).
  \item On the test set most wrong pictures of \Dp{} belong to another sound (15 of 24); pictures of on-screen sounds
        are few (4).
\end{itemize}

\subsection{Statistical comparisons}\label{sec:statres}
\begin{table}[h]
\centering\small
\caption{Paired clip bootstrap on merged TEST (2000 draws). $d$ = mean per-clip cost difference (negative = first
system cheaper); $p$ one-sided as in Section~\ref{sec:stats}; the two-sided value is about twice as large.}
\label{tab:stats}
\begin{tabular}{lrrl}
\toprule
comparison & $d$ & 95\,\% interval & one-sided $p$ \\
\midrule
\Dp{} vs B & $-0.136$ & $[-0.295,\ 0.000]$ & 0.034 \\
D vs B & $-0.159$ & $[-0.341,\ -0.023]$ & 0.0145 \\
SHIP8 vs old pipeline B0r & $-0.341$ & $[-0.636,\ -0.045]$ & 0.017 \\
first pre-registered exposure: TO1+F7F8 vs B0r & $-0.091$ & $[-0.250,\ +0.045]$ & 0.132 \\
\Dp{} vs B0r & \multicolumn{3}{l}{\emph{not yet computed} (see below)} \\
\bottomrule
\end{tabular}
\end{table}

Three cautions apply to Table~\ref{tab:stats}. (1) The $p$-values are one-sided: \Dp{} versus B is $p \approx 0.07$
two-sided, so ``\Dp{} is cheaper than B on new data'' is suggestive, not established. (2) The test set had been read
after every shipped change (Section~\ref{sec:disclosures}); none of these comparisons is a blind confirmation. (3) The
direct test of \Dp{} against the old pipeline has not been run: the per-clip rows of \Dp{} on TEST are on the cluster
(job 31708734), and the old pipeline's rows were computed with a different display rule (repeat merge 2.0\,s, no
grouping), which moves its wrong-picture count (Section~\ref{sec:mergegap}). The command to run it is in
Section~\ref{sec:repro}. The raw difference is large ($-0.500$ per clip; 3 more hits, 16 fewer wrong pictures), and
SHIP8 alone already differs from B0r with one-sided $p = 0.017$, but the number itself is not claimed here.

The step-by-step test reads of the early chain were also corrected for multiplicity (Holm over the four reads of
30 Sept): none of them was significant after correction (\fp{benchmark/gold/stack_stats.json}).

\subsection{Sensitivity}\label{sec:sensitivity}
\paragraph{On-screen weight.} With Eq.~\ref{eq:costw}, \Dp{} costs 1.887 / 1.972 / 2.056 on DEV and 2.318 / 2.364 /
2.409 on TEST at $w = 0 / 1 / 2$ (derived). For every shipped step from the old pipeline to SHIP8 the cost falls at
every $w$, so no shipping decision depends on how on-screen pictures are weighted
(\fp{benchmark/gold/visible_weight_sweep.md}).

\paragraph{Repeat-merge gap.}\label{sec:mergegap} Re-reading the same saved pictures at different merge gaps
(\fp{benchmark/gold/merge_gap_sens.py}) moves wrong-picture counts but not the ranking:
\begin{center}\small
\begin{tabular}{lcccc}
\toprule
gap (s) & SHIP8 DEV & B0r DEV & SHIP8 TEST & B0r TEST \\
\midrule
1.5 & 28 / 23 / 2.338 & 18 / 57 / 3.859 & 23 / 30 / 2.591 & 21 / 44 / 3.000 \\
2.0 & 28 / 21 / 2.282 & 18 / 51 / 3.690 & 23 / 29 / 2.568 & 21 / 40 / 2.909 \\
3.0 & 28 / 18 / 2.197 & 17 / 43 / 3.521 & 22 / 27 / 2.568 & 20 / 33 / 2.795 \\
\bottomrule
\end{tabular}
\end{center}
The old pipeline benefits more from a longer gap than SHIP8 does (its many repeated pictures are merged), which is why
\Dp{} (gap 2.5) against B0r (gap 2.0) is not a like-for-like comparison.

\paragraph{The 0.5\,s/1.0\,s window and $\beta$.} The window and $\beta$ were fixed before the detector push and not
varied for \Dp. In the September table (Section~\ref{sec:sept-table}) the gated system was the cheapest of all systems
for any $\beta$ between 0.39 and 2.56 on TEST (0.46--2.33 on DEV), so $\beta = 2$ lies inside that range.

\subsection{Where the misses and wrong pictures come from}\label{sec:errors}
The Decision Inspector records, for every needed sound, the step at which it was lost
(\fp{docs/review/inspector_trail_ready.md}):

\begin{table}[h]
\centering\small
\caption{Fate of the missed needed sounds of \Dp.}
\label{tab:misses}
\begin{tabular}{L{4.0cm}rrL{8.4cm}}
\toprule
 & DEV & TEST & where \\
\midrule
misses & 29 & 41 & \\
\quad lost at a named step & 20 & 32 & DEV: band rescue 6, gate 5, DASM vote 2, family merge 2, mirror veto 2, scene check 1, masked-weak 1, one rescue per family 1. TEST: band rescue 6, continuation veto 5, family merge 4, gate 3, mirror veto 3, DASM vote 2, BEATs extraction 2, seven other steps 1 each \\
\quad never heard by any detector & 6 & 6 & no span of the family near the onset \\
\quad timing / matching & 3 & 3 & a same-family span exists but starts outside the window \\
\bottomrule
\end{tabular}
\end{table}

Two things stand out. Most misses are not ``never heard'': the sound was a candidate and a later rule removed it,
most often because the listeners did not confirm a weak detection (band rescue). And the gate silences few needed
sounds (5 DEV, 3 TEST), but it is also the reason the on-screen wrong pictures remain (6 DEV, 4 TEST).

\subsection{How far could this pipeline go? (ceiling analysis)}\label{sec:ceiling}
The ceiling analysis replaces one stage at a time by the gold answer on the \emph{same} detections, using the saved
pictures of SHIP7 (\fp{benchmark/gold/ceiling_ship7.md}; this ran before the visibility re-check, so it counts 55
needed sounds and SHIP7 scores 25 / 27 / 2.451 there). It is the best available map for \Dp, which is SHIP7 plus later
display and witness rules.

\begin{table}[h]
\centering\small
\caption{Single-stage oracles on SHIP7 (merged DEV, 55 needed).}
\label{tab:ceiling}
\begin{tabular}{lrrrr}
\toprule
oracle (one stage made perfect) & hits & wrong (v/c/p) & cost & $\Delta$ \\
\midrule
none (SHIP7 as shipped) & 25 & 27 (9/16/2) & 2.451 & -- \\
perfect on-screen gate & 30 & 18 (0/16/2) & 1.915 & $-0.536$ \\
perfect listener, filters lifted & 34 & 24 (7/15/2) & 1.859 & $-0.592$ \\
perfect vetoes / filters & 30 & 27 (9/16/2) & 2.169 & $-0.282$ \\
perfect timing & 26 & 20 (11/7/2) & 2.197 & $-0.254$ \\
perfect family & 25 & 26 (14/10/2) & 2.423 & $-0.028$ \\
\midrule
all stages perfect (ceiling) & 43 & 3 (0/1/2) & 0.761 & $-1.690$ \\
\bottomrule
\end{tabular}
\end{table}

Applied in sequence (timing, family, gate, vetoes, listener), the gate gives the largest single step ($-0.676$). Even
with every stage perfect, 12 of 55 needed sounds stay missed: 5 are heard by no model at all (a hammer, an explosion,
a clang, two golf whacks) and 7 are too faint or too badly timed. This is the \emph{detector floor}: with today's
detectors no decision layer can find more than about four in five needed sounds.

\subsection{The value of the on-screen gate (September test table)}\label{sec:sept-table}
The last time the gated system was compared with the \emph{blind} baseline (same detector and generator, no gate) was
the September test table (amendment 21, tag \texttt{test\_final\_v33}, TEST-60; \fp{benchmark/gold/holm_test_final_v33_test_bench.json}).
It used the old pipeline's detector, not \Dp's, but the same gate model and gate rule, so it is the best available
estimate of what the gate itself is worth:

\begin{table}[h]
\centering\small
\caption{September test table: gated minus blind on TEST-60, two-sided $p$ with Holm correction.}
\label{tab:sept}
\begin{tabular}{lrrl}
\toprule
measure & gated $-$ blind & 95\,\% interval & Holm $p$ \\
\midrule
F1 (pre-registered primary) & $+0.059$ & $[-0.030,\ +0.144]$ & 0.18 (unadjusted, null) \\
precision & $+0.137$ & & 0.010 \\
wrong pictures per clip & $-0.517$ & $[-0.800,\ -0.283]$ & $<0.001$ \\
viewer cost per clip & $-0.833$ & & 0.012 \\
clean-clip accuracy & $+0.219$ & & 0.006 \\
\midrule
unseen clips: gated vs showing nothing & $-3.077$ & & $<0.001$ \\
all clips: gated vs showing nothing & $-0.233$ & & 0.585 \\
\bottomrule
\end{tabular}
\end{table}

The gate removes about half a wrong picture per clip at the price of a small loss in recall, on both halves of the
data (on DEV the precision gain was not significant after Holm, 0.080). Its value is therefore real but lies almost
entirely in what it does \emph{not} show, which is why F1 -- which weighs a miss and a wrong picture equally -- does not
move. A gate-versus-blind comparison for \Dp{} itself was not re-run.

\subsection{Other measurements}\label{sec:other-results}
\paragraph{Gate accuracy.} On the gold DEV sounds the Qwen3.8-27B gate has balanced accuracy 0.617: it silences 16 of
43 on-screen sounds (0.37) and keeps 31 of 36 needed ones (0.86) (\fp{docs/review/vlm_panel_2026-10-01/f1_audit.md}).
Its predecessor Qwen2.5-VL-7B reached 0.61. A 0.5-s shift of the sampled frames changes 6 of 79 verdicts (7.6\,\%).

\paragraph{Listener agreement as evidence.} On the held-out 415 AudioSet-Strong clips, a BEATs detection that neither
listener names is a real sound 22\,\% of the time, one named by one listener 37\,\%, and one named by both 62\,\%. This is
the basis of the witness rule, and it is why the remaining decisions sit where the models disagree at about chance.
The audio part of the best rejected idea (a listener names sounds over the whole clip, FineLAP times them, DASM
confirms) was right on 74 of 92 detections (0.80) on the same clips, against 0.38--0.39 for the raw detectors: the
extra sounds are mostly real, and they become wrong pictures in the full pipeline because the gate passes off-screen
ones whose name is wrong, or because the gold lists only the salient sounds.

\paragraph{Onset timing (v3 era).} On 255 DCASE 2025 Task 3 events, of the events BEATs detects, hysteresis plus
occlusion places the onset within 0.5\,s for 80\,\% (mean absolute error 0.35\,s; \fp{benchmark/eval_dcase_onset.json}).
Only 26\,\% of those short-clip events reached the display bar at all, which again points to detection, not timing,
as the limit.

\paragraph{Pictures.} Section~\ref{sec:pic-history}: in a blind glance test, Qwen-Image pictures were recognised 26 of
54 times against 14 of 54 for FLUX ($+0.222\ [+0.074, +0.370]$), by the author as rater.
